% =========================================================================
\section{Resumen del método y perspectivas}\label{sec:resumen}
% =========================================================================

\begin{metodo}[Receta para un punto ordinario $x_0$]
\begin{enumerate}[label=\textcolor{ambar}{\bfseries\arabic*.}, leftmargin=1.6em]
  \item Comprobar que $P(x_0)\neq0$ y estimar
        $R\ge\dist\bigl(x_0,\mathcal Z(P)\bigr)$, la distancia al cero
        \emph{complejo} de $P$ más cercano (\cref{teo:fuchs}).
  \item Proponer $y=\sum_{n\ge0}a_nz^{n}$, con $z=x-x_0$, y derivar término
        a término.
  \item Sustituir ---en la forma normal o, si $P$ es un polinomio,
        directamente en \eqref{eq:edo}---, desplazar índices para que todo
        vaya en $z^{n}$ y usar el producto de Cauchy.
  \item Igualar a cero el coeficiente de cada $z^{n}$
        (\cref{cor:unicidad}): se obtiene la recurrencia para $a_{n+2}$.
  \item Dejar libres $a_0=y(x_0)$ y $a_1=y'(x_0)$: $y_1$ corresponde a
        $(a_0,a_1)=(1,0)$ e $y_2$ a $(0,1)$; la solución general es
        $y=a_0y_1+a_1y_2$.
\end{enumerate}
\end{metodo}

La \cref{fig:flujo} recoge el mismo algoritmo en forma de diagrama, y la
\cref{tab:resumen} compara los cuatro ejemplos de la \cref{sec:ejemplos}.

\begin{table}[!htb]
\centering
\caption{Resumen de los ejemplos resueltos en $x_0=0$.}
\label{tab:resumen}
\renewcommand{\arraystretch}{2.25}
\small
\begin{tabular}{@{}l l l c l@{}}
\toprule
\textcolor{cian}{\bfseries Ecuación} & \textcolor{cian}{\bfseries Recurrencia} &
\textcolor{cian}{\bfseries Salto} & \textcolor{cian}{\bfseries $R$} &
\textcolor{cian}{\bfseries Soluciones}\\
\midrule
$y''+y=0$ & $a_{n+2}=-\dfrac{a_n}{(n+2)(n+1)}$ & $2$ & $\infty$ & $\cos x$, $\sin x$\\
$y''-xy=0$ & $a_{n+2}=\dfrac{a_{n-1}}{(n+2)(n+1)}$ & $3$ & $\infty$ & $\Ai$, $\Bi$ (no elementales)\\
$y''+xy'+y=0$ & $a_{n+2}=-\dfrac{a_n}{n+2}$ & $2$ & $\infty$ & $\eul^{-x^2/2}$, tipo Dawson\\
Legendre & $a_{n+2}=\dfrac{(n-\ell)(n+\ell+1)}{(n+2)(n+1)}\,a_n$ & $2$ & $1$ & $P_\ell$ si $\ell\in\N$; otra con $R=1$\\
\bottomrule
\end{tabular}
\end{table}

% Figura: diagrama de flujo del método
\begin{figure}[!htb]
\centering
\begin{tikzpicture}[node distance=0.55cm]
  \tikzset{paso/.style={nodo, text width=7.4cm, align=left, minimum height=0.85cm},
           num/.style={font=\sffamily\bfseries\small, text=ambar}}
  \node[nodocian, align=center, text width=5.4cm] (ini)
    {$P(x)\,y''+A(x)\,y'+B(x)\,y=0$\\[1pt]\footnotesize\color{gris} y un punto $x_0$};
  \node[diamond, aspect=2.4, draw=violeta, fill=cajafondo, line width=0.6pt,
        inner sep=1.5pt, font=\small, below=of ini] (dec) {¿$P(x_0)\neq0$?};
  \node[paso, below=0.75cm of dec] (p1)
    {\textcolor{ambar}{\bfseries 1.}\ \ $R\ge\dist\bigl(x_0,\mathcal Z(P)\bigr)$\quad(ceros complejos)};
  \node[paso, below=of p1] (p2)
    {\textcolor{ambar}{\bfseries 2.}\ \ $y=\sum a_n z^{n}$; derivar término a término};
  \node[paso, below=of p2] (p3)
    {\textcolor{ambar}{\bfseries 3.}\ \ desplazar índices y producto de Cauchy};
  \node[paso, below=of p3] (p4)
    {\textcolor{ambar}{\bfseries 4.}\ \ coeficiente de $z^{n}$ nulo $\Rightarrow$ recurrencia};
  \node[paso, below=of p4] (p5)
    {\textcolor{ambar}{\bfseries 5.}\ \ $(1,0)\mapsto y_1$, \ $(0,1)\mapsto y_2$; \ $y=a_0y_1+a_1y_2$};
  \node[nodo, draw=rojo, align=center, text width=3.6cm, right=1.6cm of dec] (sing)
    {\textcolor{rojo}{\bfseries Punto singular}\\[1pt]
     \footnotesize\color{gris} serie de Frobenius\\(Parte II)};

  \draw[flecha] (ini) -- (dec);
  \draw[flecha] (dec) -- node[right=2pt, etiqgris] {sí} (p1);
  \draw[flecha] (dec) -- node[above=2pt, etiqgris] {no} (sing);
  \foreach \a/\b in {p1/p2, p2/p3, p3/p4, p4/p5}{\draw[flecha] (\a) -- (\b);}
\end{tikzpicture}
\caption{Diagrama de flujo del método de la serie entera.}
\label{fig:flujo}
\end{figure}



\paragraph{Perspectivas}
El método de la serie entera falla cuando $P(x_0)=0$: en un punto singular
la solución puede no ser analítica (piénsese en $y=x^{r}$ con $r\notin\N$,
o en $y=\ln x$). La segunda parte del tema \cite[cap.~V]{TEMAV} estudia los
puntos \emph{singulares regulares} (\cref{obs:avance}) mediante series de
Frobenius $y=(x-x_0)^{r}\sum_n a_n(x-x_0)^n$, cuyo exponente $r$ fija la
\emph{ecuación indicial}; el caso de raíces que difieren en un entero
conduce a las \emph{series logarítmicas}. El cambio $\bar x=1/x$ permite,
además, analizar el comportamiento de las soluciones en un entorno del
infinito. Finalmente, las ecuaciones clásicas de la física ---Hermite,
Chebyshev, Legendre--- comparten la estructura de la \cref{sec:legendre}:
una recurrencia de salto dos cuya serie sólo se trunca para valores
discretos del parámetro, lo que da lugar a familias de polinomios
ortogonales y a la cuantización de magnitudes físicas
\cite{ARFKEN,RILEY}.
